% Variant A: a tree read from top to bottom; every path ends in its own result, straight below its last
% condition. The two results of d_model^2 per block share the strong blue.
\begin{tikzpicture}[x=1mm,y=1mm]
\node[title] at (0,10) {Functional dimension added by identity skips};
\node[root] (root) at (56,0) {attention or MLP blocks\\[-1.5pt]\sub{$d = \dm$}};
\node[cond] (n0) at (17,-12) {no normalization};
\node[cond] (n1) at (95.5,-12) {pre-RMSNorm};
\node[cond] (e0) at (69,-25) {$\varepsilon \to 0$};
\node[cond] (e1) at (122,-25) {$\varepsilon > 0$};
\node[cond] (h1) at (50,-38) {ReLU, ReLU$^2$, ReGLU\\[-1pt]\sub{$B(\lambda h) = \lambda^k B(h)$}};
\node[cond] (h0) at (88,-38) {attention, $\tanh$, GELU,\\SwiGLU, GeGLU};
\node[full] (r1) at (17,-63) {$+d^2$\\[-1pt]\sub{per block}};
\node[full] (r2) at (50,-63) {$+d^2$\\[-1pt]\sub{per block}};
\node[less] (r3) at (88,-63) {$+d(d-1)/2 + 1$\\[-1pt]\sub{per block}};
\node[less] (r4) at (122,-63) {$+d(d-1)/2$\\[-1pt]\sub{per block}};
\draw[edge] (root) -- (n0);
\draw[edge] (root) -- (n1);
\draw[edge] (n1) -- (e0);
\draw[edge] (n1) -- (e1);
\draw[edge] (e0) -- (h1);
\draw[edge] (e0) -- (h0);
\draw[edge] (n0) -- (r1) node[pos=0.22, pr, align=center] {proven for 1 block$^*$\\\fdref{thm:onelayer}} node[pos=0.45, sm] {simulated, up to 64 blocks};
\draw[edge] (h1) -- (r2) node[pos=0.5, sm] {simulated, 3 blocks};
\draw[edge] (h0) -- (r3) node[pos=0.5, sm] {simulated, 3 $\tanh$ blocks};
\draw[edge] (e1) -- (r4) node[pos=0.28, pr, align=center] {proven for 1 block$^*$:\\gain ${\ge}$ value, \fdref{thm:onelayer}} node[pos=0.62, sm] {simulated, up to 8 blocks};
\node[anchor=west, font=\footnotesize, text=black!60] at (0,-74) {$^*$not for SwiGLU or GeGLU};
\end{tikzpicture}
